% Source: https://artofproblemsolving.com/community/c3h2442503_100th_post_celebration
% Stable identity: aops-topic-2442503-q02
The solutions are the root $x=-a$ of $x+a=0$ and the roots of $x^2-4ax+b=0$. For there to be three different solutions, the quadratic must have two different roots $r_1,r_2$, neither equal to $-a$. By Vieta's formulas $r_1+r_2=4a$, so the sum of the three solutions is
\[4a+(-a)=3a=60,\qquad a=\boxed{20}.\]
Such equations exist: for $a=20$, $b=0$ the equation is $x(x-80)(x+20)=0$, with solutions $0,80,-20$ and sum $60$.
